% Part: lambda-calculus
% Chapter: representability
% Section: currying

\documentclass[../../../include/open-logic-section]{subfiles}

\begin{document}

\olfileid{lam}{rep}{cur} 
\olsection{التكريع}

يمثل مجرّد~$\lambd$ من الصورة~$\lambd[x][M]$ دالة ذات وسيط واحد، وهذا
قيد شديد حين نريد تعريف دالة تقبل عدة وسائط. ويمكن تحقيق ذلك بتوسيع
حساب~$\lambd$ بحيث يسمح بتكوين الأزواج والثلاثيات وما إليها؛ وفي هذه
الحالة تتوقع دالة ثلاثية الموضع~$\lambd[x][M]$، مثلًا، أن يكون وسيطها
ثلاثية. لكن الأيسر أن نفعل ذلك بواسطة \emph{التكريع}.

لنتأمل مثالًا. سنفترض لحظة أن لدينا عملية~$+$ في حساب~$\lambd$. فدالة
الجمع $2$-موضعية، أي إنها تأخذ وسيطين. لكن مجرّد~$\lambd$ لا يعطينا
إلا دوال ذات وسيط واحد؛ إذ لا يسمح النحو بتعابير مثل
$\lambd[(x,y)][(x+y)]$. غير أنه يمكننا تأمل الدالة الأحادية الموضع
$f_x(y)$ المعطاة بـ$\lambd[y][(x+y)]$، وهي تضيف~$x$ إلى وسيطها الوحيد
$y$. والحق أن هذه ليست دالة واحدة، بل عائلة من دوال مختلفة هي ``أضف
$x$''، واحدة لكل عدد~$x$. ويمكننا الآن تعريف دالة أحادية الموضع أخرى
$g$ بأنها~$\lambd[x][f_x]$. فإذا طبقت هذه الدالة على الوسيط~$x$ أعادت
$g(x)$ الدالة~$f_x$، فتكون قيمها دوال أخرى. وإذا طبقنا~$g$ على~$x$،
ثم طبقنا الناتج على~$y$، حصلنا على~$(g(x))y = f_x(y) = x+y$. وهكذا
تستطيع الدالة الأحادية الموضع~$g$ أداء العمل نفسه الذي تؤديه دالة الجمع
الثنائية الموضع. ولا يشير ``التكريع'' إلا إلى هذه الحيلة التي تحول دوال
ثنائية الموضع إلى دوال أحادية الموضع (قيمها دوال أحادية الموضع).

وفيما يأتي مثال مصوغ كما ينبغي في نحو حساب~$\lambd$. كيف نمثل الدالة
$f(x,y) = x$؟ إذا أردنا تعريف دالة تقبل وسيطين وتعيد الأول، أمكننا أن
نكتب~$\lambd[x][\lambd[y][x]]$، وهي حرفيًا دالة تقبل وسيطًا~$x$ وتعيد
الدالة~$\lambd[y][x]$. وتقبل الدالة~$\lambd[y][x]$ وسيطًا آخر~$y$، لكنها
تهمله وتعيد~$x$ دائمًا. لنر ما يحدث حين نطبق
$\lambd[x][\lambd[y][x]]$ على وسيطين:
\begin{align*}
  (\lambd[x][\lambd[y][x]])MN 
  \bredone & (\lambd[y][M])N \\
  \bredone & M
\end{align*}

وبوجه عام، لكتابة دالة معلماتها~$x_1$، \dots،~$x_n$ ويعرّفها حد~$N$،
يمكننا أن نكتب
$\lambd[x_1][\lambd[x_2][\ldots\lambd[x_n][N]]]$. وإذا طبقنا عليها
$n$ من الوسائط حصلنا على:
\begin{multline*}
  (\lambd[x_1][\lambd[x_2][\ldots\lambd[x_n][N]]]) M_1 \dots M_n \bredone\\
  \begin{aligned}
  \bredone {} & (\Subst{(\lambd[x_2][\ldots\lambd[x_n][N]])}{M_1}{x_1}) M_2
  \dots M_n\\
   \eqs {} & (\lambd[x_2][\ldots\lambd[x_n][\Subst{N}{M_1}{x_1}]]) M_2
            \dots M_n \\
  \vdots & \\
  \bredone {} &\Subst{\Subst{N}{M_1}{x_1}\ldots}{M_n}{x_n}
  \end{aligned}
\end{multline*}
ويعني السطر الأخير حرفيًا إحلال~$M_i$ محل~$x_i$ في متن تعريف الدالة،
وهو بالضبط ما نريده حين نطبق وسائط متعددة على دالة.
\end{document}
